% ============================================================
% Chapter 3: System Architecture and PCB Design
% ============================================================
\chapter{System Architecture and PCB Design}
\label{chap:architecture}

This chapter details the technical implementation of the retrofitted test-diagnosis system, translating the design principles discussed in~\Cref{sec:design_principles} into a concrete hardware architecture. The system employs a modular design across four distinct~\glspl{pcb} to enhance isolation, testability, and maintainability, reflecting best practices for complex electronic systems~\cite{jones2004}. Detailed schematics and layout diagrams, essential for understanding the implementation, are provided in~\Cref{app:schematics}.

% ------------------------------------------------------------
\section{Overall System Architecture Overview}
\label{sec:overall_architecture}

\begin{figure}[htbp]
	\centering
	\includegraphics[width=0.75\textwidth]{overall_system_block_diagram}
	\caption{High-level block diagram of the retrofitted test-diagnosis system modules (Base, Upper, STM32, LED) and key interconnections. Based on Leyer~\cite{leyer2023ppt}.}
	\label{fig:overall_system}
\end{figure}

The retrofitted system comprises four main~\gls{pcb} modules (\Cref{fig:overall_system}) designed to collaboratively replicate and enhance the functionality of the original test-diagnosis unit:

\begin{description}
	\item[Base Board (Unterplatine):] The central processing hub housing the Actel \texttt{A3P1000-PQG208}~\gls{fpga} (\Cref{fig:base_schematic_p1}, designator \texttt{Y20} in~\Cref{tab:bom}) running the emulated~\gls{mcs48} core. The ProASIC3 family was chosen for its non-volatile Flash-based configuration memory, offering inherent immunity to configuration-memory upsets (firm errors) compared to SRAM-based alternatives, and providing sufficient logic resources for the emulation task~\cite{guzman2011obsolescence}. This board integrates crucial support circuitry including power regulation (\texttt{Y1}, \texttt{Y2}), the \SI{6}{\mega\hertz} clock generation (\texttt{Q6}, \texttt{Y16}, \Cref{fig:base_schematic_p2}), \gls{nvram}~(\texttt{Y15}, \texttt{CY14B104NA}, \Cref{fig:base_schematic_p2}) for non-volatile storage of fault codes, \gls{bod} and watchdog circuits (\texttt{Y3}, \texttt{Y4}, \texttt{Q1}, \texttt{Y5}, \Cref{fig:base_schematic_p1}) for system supervision, a~\gls{jtag} interface (\texttt{X1}) for programming and debugging, and the necessary inter-board connectors (\texttt{X5}/\texttt{X6}, \texttt{X9}, \texttt{X13}). It also physically houses several level-shifter and buffer ICs (\texttt{Y7}, \texttt{Y8}, \texttt{Y17}, and associated CD4050 buffers) supporting the Upper Board interface logic. Detailed schematics and an example layout view are available in~\Cref{sec:app_base_board_schematics} (\Cref{fig:base_schematic_p1,fig:base_schematic_p2,fig:base_layout_top,fig:base_layout_bottom}). Note that the provided layout (\Cref{fig:base_layout_top,fig:base_layout_bottom}) corresponds to the inherited Rev~0.0 board which required rework (\Cref{subsec:hardware_challenges}).

	\item[Upper Board (Oberplatine):] Acts as the primary physical and electrical interface layer (\Cref{fig:oberplatine_schematic}), mounting onto the Base Board via header~\texttt{X8}. Its main functions include terminating the level-shifted signals from the Base Board, housing conditioning circuitry (e.g., debouncing) for the legacy input buttons, and incorporating its own logic ICs (\texttt{Y4}, \texttt{Y5}, \texttt{Y6}). The schematic context (\Cref{fig:oberplatine_schematic}) also shows output driving functions (like AV-Sperre) logically associated with this interface layer, although the driving ICs (\texttt{Y8}, \texttt{Y12} in the BOM) are physically located on the Base Board. Further details are in~\Cref{sec:app_upper_board_schematics} (\Cref{fig:upper_layout_top,fig:upper_layout_bottom}, depicting the corrected Rev~1.0 layout).

	\item[STM32 Diagnostics Board:] (\Cref{fig:stm32_schematic}) This separate module provides enhanced diagnostic and logging capabilities, intentionally decoupled from the safety-critical anti-skid core logic using galvanic isolation. It features an \texttt{STM32F401RET6}~\gls{mc} (\texttt{Y24}, Arm\textsuperscript{\textregistered}~Cortex\textsuperscript{\textregistered}-M4), chosen for its suitable peripheral set (\gls{spi}, \gls{sdio}, \gls{usb}, \gls{rtc}) and processing capability for the logging task, aligning with hardware used in the test fixture development~\cite{dolata2023}. It includes an~\gls{rtc} (\texttt{Y19}, \texttt{DS3231SN\#}) with battery backup (\texttt{U1}, \texttt{BS-7}) for accurate timestamping, an~\gls{sdio} interface to a MicroSD card (\texttt{J6}) for persistent logging using~\gls{fatfs}, and a~\gls{usb} Type-C interface (\texttt{X15}) for data retrieval. Crucially, communication with the Base Board~\gls{fpga} occurs via a galvanically isolated~\gls{spi} link (using optocouplers \texttt{OK5-OK8}), ensuring subsystem independence as required by safety design principles (\Cref{sec:design_principles}). Isolated power (\texttt{Y22}, \texttt{Y23}) maintains this barrier. The introduction of~\gls{usb}/\gls{sd} interfaces necessitates future cybersecurity assessment (\Cref{sec:results_discussion_risk}). See~\Cref{sec:app_stm32_board_schematics} (\Cref{fig:stm32_layout_top,fig:stm32_layout_bottom} for the Rev~1.0 layout).

	\item[LED Display Board:] Houses the two physical 7-segment~\gls{led} displays (\texttt{Disp1}, \texttt{Disp2}, \texttt{SC03-12EWA},~\cite{kingbright2020}) providing visual feedback to the user (e.g., error codes \enquote{99}, \enquote{31}, \enquote{88}), driven by~\gls{bcd}-to-7-Segment driver ICs (\texttt{Y13A}/\texttt{Y14A}, \texttt{CD4511BE}). It connects to the Base Board via direct connection using headers \texttt{X5}/\texttt{X6}. Schematics (integrated in \Cref{fig:oberplatine_schematic}) and layout details (\Cref{fig:led_layout_top,fig:led_layout_bottom} for Rev~1.0) are in~\Cref{sec:app_led_board_schematics}.
\end{description}

Physical interconnection utilizes direct board-to-board mounting (LED Board onto Base Board via \texttt{X5}/\texttt{X6}; \gls{stm32}~Board onto Base Board via \texttt{X13}) and cabled connections (Base Board to Upper Board via \texttt{X9} ribbon cable for Flash/control signals). This modular approach enhances testability and maintainability, fulfilling a key design principle (\Cref{sec:design_principles}). However, the necessity for rework on the inherited Base Board's \texttt{X5}/\texttt{X6} sockets (\Cref{subsec:hardware_challenges}) underscores the potential pitfalls when reusing previously modified hardware, reinforcing the need for thorough incoming inspection and verification. Base to LED Key connector details are mentioned in~\Cref{sec:app_connectors}.

% ------------------------------------------------------------
\section{FPGA Core Implementation Details}
\label{sec:fpga_core_details}

The core logic resides in the \texttt{A3P1000}~\gls{fpga}, implemented in~\gls{vhdl} (see supplementary files, e.g., \texttt{atop.vhd}, \texttt{t48\_core.vhd}, \texttt{t8243\_sync\_notri.vhd}, \texttt{sweep\_top.vhd}, \texttt{Logger\_top.vhd}).

\begin{description}
	\item[MCS-48 Emulation Core:] An emulation of the Intel 8048 (\gls{mcs48}) processor, specifically the mask-programmed 8048 variant identified in the legacy system, serves as the central processing element. Following Leyer's approach~\cite{leyer2023ppt}, an existing implementation (\texttt{t48\_core.vhd}) from the OpenCores repository\footnote{OpenCores website: \url{https://opencores.org/}. Validation of the specific core used relies on project testing.} was integrated \enquote{as-is}. This core executes the original \SI{2}{\kilo B} firmware binary (stored in the external Flash~\gls{rom}), preserving the legacy control logic and ensuring behavioural compatibility at the instruction level. The primary verification challenge lies in ensuring this emulation is functionally and temporally equivalent to the original hardware within the target application's constraints (\Cref{subsec:equivalence_strategies,subsec:fpga_certification}).
	\item[Memory Interfaces:] The~\gls{fpga} implements the necessary logic to interface directly with the external memory components located on the Base Board: an \texttt{S29AL016J} Flash~\gls{rom} (\texttt{Y21}) storing the executable firmware, and a \texttt{CY14B104NA}~\gls{nvram} (\texttt{Y15}) used for persistent storage of fault codes and potentially other operational data (\Cref{fig:base_schematic_p2}). The~\gls{fpga} generates the required address, data, and control signals (Chip Enable, Output Enable, Write Enable) according to the timing requirements of these memory devices and the access patterns dictated by the emulated~\gls{mcs48} core. Ensuring correct timing on these interfaces was a key part of Layer~1 verification (\Cref{sec:fpga_interface_testing}).
	\item[Peripheral Emulation and Interfaces:] The~\gls{fpga} logic replicates or interfaces with the essential peripherals of the legacy test-diagnosis board:
	      \begin{itemize}
		      \item \textit{7-Segment Display Driving Logic:} Generates the~\gls{bcd} signals and the latch enable (\texttt{FPGA\_7SEG\_LE}) required by the external \texttt{CD4511BE} display drivers located on the LED Board (\Cref{fig:oberplatine_schematic}).
		      \item \textit{Input Button Handling:} Monitors the state of the \texttt{T0} and \texttt{T1} inputs (connected via the Upper Board to legacy buttons STW and Lösch) and generates interrupts for the emulated~\gls{mcs48} core as appropriate. It also reads the state of \texttt{P10}/\texttt{P11} (Test, Tür/V) (\Cref{fig:base_schematic_p2}).
		      \item \textit{Frequency Sweep Generation:} An internal~\gls{vhdl} state machine (\texttt{sweep\allowbreak\_top.vhd}) replicates the behaviour of the original analog~\gls{vco}-based circuit, generating the required frequency sweep patterns (\SI{850}{\hertz} to \SI{1550}{\hertz} and vice-versa) for the self-test function based on control signals from the emulated core. Correct timing and frequency were verified (\Cref{subsec:fpga_interface_results}).
		      \item \textit{\gls{spi} Master Interface:} Implements an~\gls{spi} master protocol (\texttt{Logger\_top.vhd}) to transmit diagnostic data (error states, button presses) to the~\gls{stm32} diagnostics board slave interface (\Cref{sec:stm32_details}).
		      \item \textit{Watchdog/Reset Interaction:} Manages the interaction with the external hardware watchdog (\Cref{fig:base_schematic_p1}), receiving the master reset signal (\texttt{!FPGA\_8048\_RESET}) and generating the necessary interrupt signal (\texttt{!INT}) to periodically reset the external watchdog under normal operation, preventing a system reset.
	      \end{itemize}
	\item[Clocking:] An external \SI{6}{\mega\hertz} oscillator (\texttt{Q6} on \Cref{fig:base_schematic_p2}) provides the primary clock source for the~\gls{fpga} and the emulated~\gls{mcs48} core, ensuring timing consistent with the original system design.
\end{description}
The~\gls{vhdl} implementation aims for functional and timing equivalence, verified through the multi-layered strategy described in~\Cref{chap:testing}.

% ------------------------------------------------------------
\section{STM32 Diagnostics Subsystem Details}
\label{sec:stm32_details}

The dedicated diagnostics board (\Cref{fig:stm32_schematic}) provides enhanced monitoring capabilities, designed to be functionally independent from the safety-critical~\gls{fpga} core. Its purpose is solely for logging and diagnostics, not influencing the primary anti-skid safety function.

\begin{description}
	\item[Microcontroller:] An \texttt{STM32F401RET6} (\texttt{Y24}, Arm\textsuperscript{\textregistered}~Cortex\textsuperscript{\textregistered}-M4) provides the processing power and peripheral set for the diagnostic functions. This microcontroller was chosen for its suitable performance, integrated peripherals (\gls{spi}, \gls{sdio}, \gls{usb}, \gls{rtc}), availability, and familiar development environment (STM32CubeIDE), and aligns with the hardware used in the related test fixture project~\cite{dolata2023} (which used an \texttt{STM32F401RE} Nucleo board).
	\item[Galvanically Isolated SPI Slave:] Optocouplers (\texttt{OK5-OK8} on \Cref{fig:stm32_schematic}) provide galvanic isolation for the~\gls{spi} communication link (MISO, MOSI, SCK, SS) from the Base Board~\gls{fpga}. This crucial isolation prevents electrical faults or failures within the non-safety-critical diagnostics board from propagating to the safety-related~\gls{fpga} domain, a key safety design principle (\Cref{sec:design_principles}). The~\gls{stm32} acts as the~\gls{spi} slave, receiving data packets containing error and status information from the~\gls{fpga} master. The C firmware module \texttt{msg\_receiver.c} (derived from~\cite{dolata2023}) handles the reception and interpretation of these packets.
	\item[Data Logging:] An~\gls{sdio} interface connects to a standard MicroSD card (\texttt{J6}), enabling persistent storage of timestamped event logs. The~\gls{stm32} firmware utilizes the~\gls{fatfs} library (\texttt{ff.c}) to manage the file system on the~\gls{sd}~card, allowing logs to be written in a standard, accessible format (e.g., text files).
	\item[Real-Time Clock (RTC):] An onboard~\gls{rtc} (\texttt{Y19}, \texttt{DS3231SN\#}) with battery backup (via holder \texttt{U1}, \texttt{BS-7}) provides accurate date and timestamps for logged events, even across power cycles. The time can be initially set via a \texttt{time.txt} file placed on the~\gls{sd}~card, as implemented by Leyer~\cite{leyer2023}.
	\item[USB Interface:] A~\gls{usb} Type-C connector (\texttt{X15}) allows the~\gls{sd}~card to be accessed as a standard Mass Storage Device when connected to a PC, facilitating easy retrieval of log files without needing to remove the card. This interface, while convenient, introduces potential cybersecurity considerations (\Cref{sec:results_discussion_risk}) that would need assessment before operational deployment.
	\item[Isolated Power Supply:] Onboard isolated DC/DC converters (\texttt{Y22}, \texttt{Y23}) create the necessary power domains for the~\gls{stm32} and its peripherals, fed via isolated connections (\texttt{X12}/\texttt{X7}) from the Base Board, maintaining the galvanic isolation barrier.
	\item[Status Indication:] An onboard~\gls{led} (\texttt{LED3}) provides basic status feedback for the diagnostics subsystem (e.g., indicating power or logging activity).
\end{description}
The C firmware, developed using STM32CubeIDE (see \texttt{GROUP\_C\allowbreak\_Retrofit\allowbreak\_STM32\allowbreak\_Firmware} in supplementary materials, e.g., \texttt{main.c}, \texttt{data\allowbreak\_map.c}, \texttt{state\allowbreak\_machine.c}, \texttt{msg\allowbreak\_receiver.c} based on the \texttt{STM32F401RE} configuration in \texttt{esp-comunication.ioc}), manages the~\gls{spi} slave communication, \gls{rtc}~timekeeping, \gls{sdio}~card access via~\gls{fatfs}, and the~\gls{usb} Mass Storage interface.

% ------------------------------------------------------------
\section{PCB Design and Implementation}
\label{sec:pcb_design_impl}

The physical design of the Rev~1.0~\glspl{pcb} focused on modularity, reliability, and compliance, specifically addressing the~\gls{dfm} issues identified in the prior prototype (\Cref{subsec:hardware_challenges}).

\begin{description}
	\item[Modularity Benefits:] Separation onto four distinct boards (Base, Upper, STM32, LED) enhances electrical isolation between functional blocks (especially safety vs.~non-safety) and significantly improves testability and maintainability during development and potentially in service. This directly addresses the modularity principle from \Cref{sec:design_principles}.
	\item[Layout Practices:] Standard layout principles, as outlined by Jones~\cite{jones2004}, were applied, informed by internal guidelines (e.g., project guidelines, adapted). This included the use of multi-layer boards (4-layer) allowing for dedicated ground planes (\Cref{fig:base_layout_top}~shows an example top copper layer, ground planes are typically internal) to ensure signal integrity and low-impedance return paths. Extensive decoupling capacitors were placed close to~\gls{ic} power pins, and trace routing followed best practices (e.g., avoiding sharp angles, using 45-degree bends, appropriate width based on current/impedance, see~\cite[p.~7 \& 10]{jones2004}).
	\item[\gls{emc}/\gls{emi} Mitigation:] Techniques targeting \gls{en50155}:\penalty0 2021/\penalty0 \gls{en50121} requirements~\cite{midya2013emc, en50121_3_2_2016} were incorporated. These include the use of ground planes, widespread decoupling, physical separation of potentially noisy digital sections from sensitive analog inputs (though less critical on this specific board), and the crucial galvanic isolation of the~\gls{stm32} subsystem via optocouplers.
	\item[Addressing Rework and DFM Issues:] The Rev~1.0~\gls{pcb} layouts explicitly corrected the footprint and connectivity errors identified in Leyer's prototype~\cite{leyer2023}. Challenges associated with the inherited Base Board Rev~0.0 (specifically the \texttt{X5}/\texttt{X6} connectors) were mitigated through careful cleaning and verification before assembly. The final assembly process was guided by the corrected schematics (\Cref{app:schematics}) and the verified~\gls{bom} (\Cref{app:bom_drc}).
	\item[Connectors:] Standard, readily available headers and sockets (e.g., pin headers, ribbon cable connectors, TE~AMPMODU System~50 cited in~\Cref{sec:app_connectors}) were used for inter-board connections (e.g., \texttt{X9}, \texttt{X13}) to ensure component availability and ease of assembly/disassembly.
\end{description}
The resulting~\glspl{pcb}, documented in~\Cref{app:schematics}, form a corrected and robust hardware implementation reflecting established electronic design principles for reliable operation.

% ------------------------------------------------------------
\section{Redundancy and Fail-Safe Features}
\label{sec:redundancy_failsafe}

Several design features were incorporated to enhance system reliability and contribute to fail-safe behaviour, aligning with functional safety principles (\Cref{sec:design_principles}).

\begin{description}
	\item[External Hardware Watchdog Timer:] An independent hardware watchdog circuit (\texttt{Y4}, a CD40106 oscillator, and \texttt{Y5}, a CD4518 counter, see~\Cref{fig:base_schematic_p1}) monitors the activity of the emulated~\gls{mcs48} core within the~\gls{fpga} via the toggling~\texttt{!INT} signal. If the~\gls{fpga} fails to toggle this signal within the expected timeframe (indicating a potential firmware hang or hardware fault), the external watchdog asserts the main reset line (\texttt{!FPGA\_8048\_RESET}), forcing the system into a defined reset state. This provides a layer of supervision independent of the core processor emulation itself. Its correct timing interaction was verified in Layer~1 testing (\Cref{subsec:fpga_interface_results}).
	\item[Brown-Out Detection (BOD):] A dedicated voltage supervisor~\gls{ic} (\texttt{TPS3307-18}, \texttt{Y3} on~\Cref{fig:base_schematic_p1}) monitors critical power supply rails (\SI{5}{\volt} and \SI{3.3}{\volt} via dividers R102--R105, and \SI{1.5}{\volt} directly). If any of these rails drop below a safe operating threshold, the~\gls{bod} circuit asserts the system reset line, preventing unpredictable behaviour during power dips or brown-outs. This ensures a defined state during power instability.
	\item[Power Supply Supervision:] The use of dedicated linear and switching regulators (\texttt{Y1}, \texttt{Y2}) ensures stable power rails are provided to the sensitive electronic components (\gls{fpga}, \gls{stm32}, memories) from the potentially variable input supply voltage. This basic supervision contributes to reliable operation.
	\item[Galvanic Isolation:] As previously detailed (\Cref{sec:stm32_details}), the opto-isolated~\gls{spi} interface between the~\gls{fpga} (safety domain) and the~\gls{stm32} (non-safety domain) acts as a fail-safe barrier, preventing electrical faults from propagating between these subsystems. This is a crucial element of the design's safety architecture.
	\item[Built-In Self-Tests (Legacy):] While not a feature of the new design itself, the ability of the retrofitted system to execute the original system's self-tests (initiated via the \enquote{Test} button) is a crucial element for demonstrating functional equivalence. These legacy tests, executed by the emulated~\gls{mcs48} core within the~\gls{fpga}, provide a baseline check of core functionality, including the frequency sweep generator and display outputs. Verification of these self-tests formed a key part of the Layer~2 system testing (\Cref{subsec:system_scenario_results}). They contribute to fault detection, though they do not constitute a modern, comprehensive~\gls{bist} as described in literature~\cite{latifshabgahi2009redundancy}.
\end{description}
These hardware features contribute significantly to the overall dependability and safety argument for the retrofitted system.